\documentclass[10pt,a4paper]{amsart}
\usepackage[margin=19mm]{geometry}
\usepackage{amsmath,amssymb,mathtools}
\usepackage[scr=rsfs,cal=euler]{mathalfa}
\usepackage{xfrac}
\usepackage[unicode,hidelinks,bookmarks=false]{hyperref}
\newcommand{\ie}{i.e.\ }
\newcommand{\cf}{cf.\ }
\newcommand{\resp}{respectively\ }
\newcommand{\Subsection}{\subsection}
\providecommand{\nobreakdash}{\nobreak\hbox{-}}
\setlength{\emergencystretch}{2em}
\multlinegap0pt
\usepackage{amssymb}
\usepackage{xcolor}
\newtheorem{lem}{Lemma}
\newtheorem{claim}{Claim}
\newtheorem{proofclaim}{Proofclaim}
\DeclareMathOperator{\Homeo}{Homeo}
\title{Filtered Boolean powers of finite simple non-abelian Mal'cev algebras}
\author{Peter Mayr, Nik Ru{\v s}kuc}
\date{Source: arXiv:2404.17322v3 (CC BY 4.0)}
\begin{document}
\maketitle

\noindent\emph{Source and licence.} Filtered Boolean powers of finite simple non-abelian Mal'cev algebras. Version arXiv:2404.17322v3 (CC BY 4.0), distributed by arXiv under the Creative Commons Attribution 4.0 International licence, http://creativecommons.org/licenses/by/4.0/, as recorded in the arXiv licence metadata for this exact version and shown on the abstract page https://arxiv.org/abs/2404.17322v3. Full licence notice: \texttt{licenses/arxiv-2404.17322-CC-BY-4.0.txt}.\par\smallskip

\noindent\textit{Editorial scope.} Counted unit: the article's claim cl:GbGcGb (source0.tex lines 2041--2045). The article's complete original proof of it is delivered with it (source0.tex lines 2048--2158, one \texttt{proof} environment of 111 lines). No mathematics is written, repaired or re-derived: the statement and the proof are the article's own lines, in the article's own notation, and no retained line is substituted or re-flowed.  Retained uncounted as the source-local context the proof needs: lines 1749--1752.  Nothing was omitted from any retained span, so the package carries no omission record.  No retained line contains a citation, so the wrapper carries no bibliography and no bibliography entry.  No retained line contains a figure environment, an image-inclusion command or drawing code, so no image is delivered and the package declares no asset dependency.  Editorial formatting only, in addition: an AMS \texttt{amsart} preamble that restates the article's own declarations and macros used by the retained lines, namely the environments \texttt{lem}, \texttt{claim}, \texttt{proofclaim} on separate counters, together with the standard packages those lines need.  The retained environments print in this document as Lemma 1, Claim 1, Proofclaim 1 in order of appearance, whereas the article prints the same items with the numbering of its own full document.  The complete native source of the article as distributed in the arXiv e-print for this exact version is bundled unchanged as \texttt{sources/158208/source0.tex}.
\begin{lem} \label{lem:types}
 For $I,I'\subseteq [n]$ with $I\cup I'=[n]$, the set of clopens in $X^\circ$ of type $(I,I')$ is an orbit under
 $(\Homeo 2^\omega)_{\{x_1,\dots,x_n\}}$.  
\end{lem}

\begin{claim} \label{cl:GbGcGb}
Let $b,c,d$ be pairwise disjoint good clopens such that $X^\circ=b\cup c\cup d$ and let $\sigma\in G$ such that  
 $d\setminus\sigma^{-1}(b)$ is good.
Then $\sigma\in G_bG_cG_b$ (pointwise stabilizers).
\end{claim}

\begin{proofclaim}
We aim to define three homeomorphisms $\sigma_1,\sigma_3\in G_b$, $\sigma_2\in G_c$ such that
$\sigma=\sigma_3\sigma_2\sigma_1$,

 Since $(c\cup d)\cap \sigma^{-1}(b)$ is clopen in $X^\circ$ and $d$ is good,
 we have $\tau_1\in G$ such that
\begin{equation} \label{eq:Gb1}
 f := \tau_1((c\cup d)\cap \sigma^{-1}(b)) \text{ is contained in } d \text{ and } d\setminus f \text{ is good}.
\end{equation}
By assumption, the clopen $d\setminus \sigma^{-1}(b)$ is good; in particular, $x_1,\dots, x_n$ belong to its closure.
 Hence the set $(c\cup d)\setminus \sigma^{-1}(b)$ also has all $x_i$ as limit points.
 The complement of this set contains the good clopen $b$, and hence also has $x_1,\dots, x_n$ among its
 limit points.
 We conclude that  $(c\cup d)\setminus \sigma^{-1}(b)$ is good.
 An analogous argument shows that $(c\cup d)\setminus f\in \mathfrak{K}$.
By Lemma \ref{lem:types} there exists a homeomorphism
\begin{equation*}
\tau_2 \colon (c\cup d)\setminus\sigma^{-1}(b)\rightarrow (c\cup d)\setminus f.
\end{equation*}
We can now define the first of our three target homeomorphisms on $X$:
\begin{equation}
\label{eq:Gb6}
\sigma_1(x) := 
\begin{cases}
x & \text{if } x\in b \cup \{x_1,\dots,x_n\}, \\
\tau_1(x) & \text{if } x\in (c\cup d)\cap \sigma^{-1}(b),\\
\tau_2(x) & \text{if } x\in (c\cup d)\cap \sigma^{-1}(c\cup d).
\end{cases}
\end{equation}
To see that $\sigma\in G$, observe that it is defined piecewise by homeomorphisms on three clopens that partition $X^\circ$,
and that the images of those three clopens, namely $b$, $f$ and $(c\cup d)\setminus f$,
also partition $X^\circ$. That $\sigma\in G_b$ is immediate from \eqref{eq:Gb6}.

Next the clopen
\[
(b\cap \sigma^{-1}(c\cup d))\cup (d\setminus f)
\]
is good: indeed, $x_1,\dots,x_n$ are in its closure
because $d\setminus f$ is good by \eqref{eq:Gb1}, while its complement also has $x_1,\dots,x_n$ in its closure
because it contains $c\in\mathfrak{K}$.
Pick a homeomorphism
\begin{equation*}
\tau_3\colon \bigl(b\cap \sigma^{-1}(c\cup d)\bigr)\cup (d\setminus f)\rightarrow d.
\end{equation*}
 Define our second target homeomorphism
 by
 \begin{equation}
 \label{eq:Gb11}
 \sigma_2(x):= 
 \begin{cases}
 x & \text{if } x\in c \cup \{x_1,\dots,x_n\},\\
 \sigma\tau_1^{-1}(x) & \text{if } x\in f,\\
 \sigma(x) & \text{if } x\in b\cap \sigma^{-1}(b),\\
 \tau_3(x) & \text{if } x\in \bigl(b\cap \sigma^{-1}(c\cup d)\bigr)\cup (d\setminus f).
 \end{cases}
 \end{equation}
 Then $\sigma_2\in G_c$ since it is defined piecewise by homeomorphisms on clopens partitioning $X^\circ$ with
 images $c, \sigma(c\cup d)\cap b, \sigma(b)\cap b, d$ also partitioning $X^\circ$.
 
 Finally, define
 \begin{equation}
 \label{eq:Gb13}
 \sigma_3(x) := 
 \begin{cases}
 x & \text{if } x\in b \cup \{x_1,\dots,x_n\},\\
 \sigma\tau_2^{-1}(x) & \text{if } x\in c,\\
 \sigma\tau_3^{-1}(x) & \text{if } x\in \tau_3(b\cap \sigma^{-1}(c\cup d)),\\
 \sigma\tau_2^{-1}\tau_3^{-1}(x) & \text{if } x\in \tau_3(d\setminus f).
 \end{cases}
 \end{equation}
 That $\sigma_3\in G_b$ follows as usual since the images of its constituting partial homeomorphisms
 $b, \sigma\tau_2^{-1}(c), \sigma(b)\cap (c\cup d), \sigma\tau_2^{-1}(d\setminus f)$ partition $X^\circ$.
 
 Now we verify that $\sigma=\sigma_3\sigma_2\sigma_1$ using the definitions \eqref{eq:Gb6}, \eqref{eq:Gb11}, \eqref{eq:Gb13}
 of $\sigma_1,\sigma_2,\sigma_3$ respectively.
 For $x\in X^\circ$
 \begin{align*}
 &\sigma_3\sigma_2\sigma_1(x)
 \\
 =&
 \sigma_3\sigma_2\left( 
 \begin{array}{ll}
 x & \text{if } x\in b\\
 \tau_1(x) & \text{if } x\in (c\cup d)\cap \sigma^{-1}(b)\\
 \tau_2(x) & \text{if } x\in (c\cup d)\cap \sigma^{-1}(c\cup d)
 \end{array}
 \right)
 \\
 =&
 \sigma_3\left(\begin{array}{ll}
 \sigma(x) & \text{if } x\in b\cap \sigma^{-1}(b)\\
 \tau_3(x) & \text{if } x\in b\cap \sigma^{-1}(c\cup d)\\
 \sigma(x) & \text{if } x\in (c\cup d)\cap \sigma^{-1}(b)\\
 \tau_2(x) & \text{if } x\in (c\cup d)\cap \sigma^{-1}(c\cup d) \cap \tau_2^{-1}(c)\\
 \tau_3\tau_2(x) &\text{if } x\in (c\cup d)\cap \sigma^{-1}(c\cup d)\cap \tau_2^{-1}(d\setminus f)
 \end{array}\right)
 \\
 =&
 \begin{cases}
 \sigma(x) & \text{if } x\in b\cap \sigma^{-1}(b)\\
 \sigma\tau_3^{-1}\tau_3(x) & \text{if } x\in b\cap \sigma^{-1}(c\cup d)\\
 \sigma(x) & \text{if } x\in (c\cup d)\cap \sigma^{-1}(b)\\
 \sigma\tau_2^{-1}\tau_2(x) & \text{if } x\in (c\cup d)\cap \sigma^{-1}(c\cup d)\cap \tau_2^{-1}(c)\\
 \sigma\tau_2^{-1}\tau_3^{-1}\tau_3\tau_2(x) & \text{if } 
 x\in (c\cup d)\cap \sigma^{-1}(c\cup d)\cap \tau_2^{-1}(d\setminus f)
 \end{cases}
 \\
 =& \sigma(x).
 \end{align*}
Thus Claim~\ref{cl:GbGcGb} is proved.
\end{proofclaim}

\end{document}
